%
% Einfache LaTeX-Vorlage f�r Arbeiten am Lehrstuhl Kranzlm�ller / MNM-Team
% - optimiert f�r die Arbeit mit g�ngigen LaTeX-Editoren
% - funktioniert ohne Makefile und Anpassungen der LaTeX-Verzeichnisstruktur
% - verwendet Komaskript f�r ein (nach europ�ischen Gepflogenheiten) sch�neres Layout
% 
% v1, 2007 (Michael Brenner)
% Diese Version: v1.1, 2012 (Michael Brenner)
% Diese Version: v1.2, 2017 (Michael Brenner)
% 


\documentclass[bibliography=totoc,listof=totoc,BCOR=5mm,DIV=12]{scrbook} % Rand f�r Bindung: 5mm / falls Index verwendet, erg�nze "index=totoc" zu den Optionen 
\usepackage{ngerman} % Trennung nach neuer deutscher Rechtschreibung, deutsche Sonderzeichen, z.B. \glqq und \grqq f�r deutsche Anf�hrungszeichen 
\usepackage[latin1]{inputenc} % Umlaute im Text
\usepackage{graphicx} % Einf�gen von Grafiken  - f�r PDF-Latex: .pdf und .png (.jpg m�glich, sollte aber vermieden werden)
\usepackage{url}           % URL's (z.B. in Literatur) sch�ner formatieren
\usepackage{hyperref} % sorgt f�r f�r Hyperlinks in PDF-Dokumenten
\usepackage{color}

\graphicspath{{./Bilder/}}

%
% der Befehl \hypenation versteht keine Sonderzeichen, also weder �
% noch "a noch \"a. W�rter die derartige Zeichen enthalten m�ssen
% direkt im Text getrennt werden, z.B. W�r\-ter
%
\hyphenation{Ma-nage-ment}
\hyphenation{Ma-nage-ment-agent}
\hyphenation{Ma-nage-ment-agent-en}
\hyphenation{Ma-nage-ment-ar-chi-tek-tur}
\hyphenation{Ma-nage-ment-ar-chi-tek-tu-ren}
\hyphenation{Ma-nage-ment-an-wen-dung}
\hyphenation{Ma-nage-ment-an-wen-dung-en}
\hyphenation{Ma-nage-ment-an-for-der-ung}
\hyphenation{Ma-nage-ment-funk-ti-on}
\hyphenation{Ma-nage-ment-funk-ti-onen}
\hyphenation{Ma-nage-ment-kon-zep-te}
\hyphenation{Ma-nage-ment-res-source}
\hyphenation{Ma-nage-ment-in-for-ma-ti-on}
\hyphenation{Ma-nage-ment-res-sour-cen}
\hyphenation{ma-nage-ment-re-le-vante}
\hyphenation{ma-nage-ment-sy-stem}
\hyphenation{ma-nage-ment-sy-steme}
\hyphenation{Ma-nage-ment-in-stru-men-tie-rung}
\hyphenation{Ma-nage-ment-platt-form}
\hyphenation{Sys-te-men}
\hyphenation{Sys-tem-um-ge-bun-gen}
\hyphenation{Sys-tem-ma-nage-ment}
\hyphenation{DHCP}
\hyphenation{Ma-nage-ment-diszi-plinen}
\hyphenation{System-management-architekturen}
\hyphenation{Verwendungs-nachweise}
\hyphenation{Video-einricht-ungen}
\hyphenation{Res-source}
\hyphenation{Res-sourcen}
\hyphenation{Grund-anwendung}
\hyphenation{Grund-anwendungen}
\hyphenation{Basis-anwendung}
\hyphenation{Core}
\hyphenation{Kom-mu-ni-ka-ti-on}
\hyphenation{De-sign-ent-schei-dung}
\hyphenation{Sprung-ad-res-sen}
\hyphenation{Klas-si-fi-ka-ti-on}
\hyphenation{Schreib-recht}
\hyphenation{Be-nut-zer-zer-ti-fi-kat}
\hyphenation{Bau-stein-ent-wi-ckler}
\hyphenation{ad-mi-ni-stra-ti-ve}

 % in dieses File kommen W�rter die Latex nicht richtig trennt

\begin{document}

% ---------------------------------------------------------------
\frontmatter % Titelbl�tter und Erkl�rung jeweils spezifisch f�r die jeweilige Uni einbinden
    %%%%%%%%%%%%%%%%%%%%%%%%%%%%%%%
% erste Seite

\thispagestyle{empty}

\begin{center}

\vspace*{-2cm}

{\Huge INSTITUT F�R INFORMATIK\\[1mm]}
DER LUDWIG--MAXIMILIANS--UNIVERSIT�T M�NCHEN\\

\vspace*{1cm}

\includegraphics[width=0.3\textwidth]{lmu_siegel}

\vspace*{2cm}

{\Large \textbf{Master's Thesis}}\\ % oder Fortgeschrittenenpraktikum, Master's Thesis, Bachelorarbeit etc.

\vspace{2.0cm}
{\Huge \textbf{C++ Graph Concepts for}}\\
\vspace*{3mm}
{\Huge \textbf{Partitioned Global Address Space}}\\
\vspace*{20mm}

{\LARGE Stefan Effenberger} % Name des Autors

\vspace{3cm}
\end{center}

\newpage

%%%%%%%%%%%%%%%%%%%%%%%%%%%%%%%
% zweite Seite

\thispagestyle{empty}
\cleardoublepage

%%%%%%%%%%%%%%%%%%%%%%%%%%%%%%%
% dritte Seite (Kopie der ersten)

\thispagestyle{empty}

\begin{center}

\vspace*{-2cm}

{\Huge INSTITUT F�R INFORMATIK\\[1mm]}
DER LUDWIG--MAXIMILIANS--UNIVERSIT�T M�NCHEN\\

\vspace*{1cm}

\includegraphics[width=0.3\textwidth]{lmu_siegel}

\vspace*{2cm}

{\Large \textbf{Master's Thesis}}\\ % oder Fortgeschrittenenpraktikum, SEP etc.

\vspace{2.0cm}
{\Huge \textbf{C++ Graph Concepts for}}\\
\vspace*{3mm}
{\Huge \textbf{Partitioned Global Address Space}}\\
\vspace*{20mm}

{\LARGE Stefan Effenberger} % Name des Autors
\vspace{2cm}

\parbox{1cm}{
\begin{large}
\begin{tabbing}
Aufgabensteller: \hspace{.5cm} \=Prof. Dr. Dieter Kranzlm�ller\\[2mm]
Betreuer:
\>Tobias Fuchs\\[5mm]
Abgabetermin: \> \textcolor{red}{ADD DATE}\\
\end{tabbing}
\end{large}}\\
\vspace{5mm}

\end{center}
 % Titelbl�tter LMU - auskommentieren falls TUM-Arbeit
%    % Richtlinien, siehe http://wwwpa.in.tum.de/generell/Abschlussarbeitsform.html
%
%%%%%%%%%%%%%%%%%%%%%%%%%%%%%%%


% Deckblatt

\thispagestyle{empty}

\begin{center}
    \includegraphics[width=3cm]{tum-logo}\\
    \vspace{.5cm}
% "Technische Universit�t M�nchen" oder alternativ das Logo der TUM
    {\Large \sc Technische Universit�t M�nchen}\\

    \vspace{1cm}
% "Fakult�t f�r Informatik"
    {\Huge \sc Fakult�t f�r Informatik\\[1mm]}


    \vspace{2cm}
% Diplomarbeit | Master's Thesis | Bachelorarbeit in Informatik | Wirtschaftsinformatik |
    {\Large \textbf{Masterarbeit in Informatik}}\\
% Thema bzw. Titel der Arbeit  (In der Sprache, in der die Arbeit verfasst wurde)
    \vspace{2.0cm}
    {\Huge \textbf{Ein Lorem-Rahmenwerk}}\\ % bei langen Titeln ggf. Schriftgr��e auf \huge herunter setzen
    \vspace*{3mm}
    {\Huge \textbf{f�r Ipsum-Systeme}}\\
    \vspace*{3mm}
    {\Huge \textbf{-- ein Dolor-Ansatz}}\\
    \vspace{1.5cm}
% Vorname und Nachname des Bearbeiters/ der Bearbeiterin
    Vorname Nachname

    \vspace{5cm} % ggf. je nach Zeilenzahl und Schriftgr��e des Titels anpassen
    \includegraphics[width=2.4cm]{tum-info-logo}
\end{center}

\newpage

%%%%%%%%%%%%%%%%%%%%%%%%%%%%%%%
% R�ckseite Deckblatt

\thispagestyle{empty}
\cleardoublepage

%%%%%%%%%%%%%%%%%%%%%%%%%%%%%%%
% Erste Seite (Titelblatt)

\thispagestyle{empty}

\begin{center}

    \includegraphics[width=3cm]{tum-logo}\\
    \vspace{.5cm}
    {\Large \sc Technische Universit�t M�nchen}\\


    \vspace{.5cm}

    {\huge \sc Fakult�t f�r Informatik\\[1mm]}


    \vspace{1cm}

    {\Large \textbf{Diplomarbeit in Informatik}}\\ % oder SEP etc.

% Thema bzw. Titel der Arbeit  (In der Sprache, in der die Arbeit verfasst wurde)
    \vspace{1.5cm}
    {\huge \textbf{Ein Lorem-Rahmenwerk}}\\ % bei langen Titeln ggf. Schriftgr��e herunter setzen
    \vspace*{3mm}
    {\huge \textbf{f�r Ipsum-Systeme}}\\
    \vspace*{3mm}
    {\huge \textbf{-- ein Dolor-Ansatz}}\\

% die englische bzw. deutsche Entsprechung des Titels
    \vspace{1cm}
    {\huge \textbf{A Lorem Framework}}\\ % bei langen Titeln ggf. Schriftgr��e herunter setzen
    \vspace*{3mm}
    {\huge \textbf{for Ipsum Systems}}\\
    \vspace*{3mm}
    {\huge \textbf{-- a Dolor Approach}}\\
    \vspace{1cm}

    \parbox{1cm}{
      \begin{large}
        \begin{tabbing}
          Bearbeiter: \hspace{1.5cm}
            \=Vorname Nachname\\[2mm]
    Aufgabensteller: \>Prof. Dr. Dieter Kranzlm�ller\\[2mm]
    Betreuer: \>MNM-Team-Betreuer 1\\ % alphabetische Reihenfolge (Nachname)
    \>MNM-Team-Betreuer 2\\
    \>Externer Betreuer 1 (Firma)\\[5mm]
    Abgabedatum: \> 7. Juli 2077\\
        \end{tabbing}
      \end{large}
    }\\

    \vspace{.3cm}

    \includegraphics[width=2.4cm]{tum-info-logo}

\end{center}
 % Titelbl�tter TUM - auskommentiert lassen falls LMU-Arbeit
    \thispagestyle{empty}
    \cleardoublepage
    %
% LaTeX-Rahmen f�r Arbeiten am Lehrstuhl Hegering
%
% Harald Roelle, 2001, 2002
%
% basierend auf Arbeiten von Helmut Reiser, Boris Gruschke und Stephen Heilbronner
%

\newpage

\thispagestyle{empty}

\begin{large}

\vspace*{2cm}

\noindent
Hiermit versichere ich, dass ich die vorliegende Masterarbeit
selbst�ndig verfasst und keine anderen als die angegebenen Quellen
und Hilfsmittel verwendet habe.

\vspace{2cm}

\noindent
M�nchen, den \textcolor{red}{ADD DATE}

\vspace{3cm}

\hspace*{7cm}%
\dotfill\\
\hspace*{8.5cm}%
\textit{(Unterschrift des Kandidaten)}

\end{large}
 % Erkl�rung (Arbeit selbstst�ndig verfasst) - auskommentieren falls TUM-Arbeit
%    \begin{large}

\vspace*{2cm}
\noindent
Ich versichere, dass ich diese Masterarbeit % (bzw. Master's Thesis)
selbst�ndig verfasst und nur die angegebenen Quellen und Hilfsmittel verwendet habe.

\vspace{2cm}

\noindent
M�nchen, den 7. Juli 2077

\vspace{3cm}

\hspace*{7cm}%
\dotfill\\
\hspace*{8.5cm}%
\textit{(Unterschrift des Kandidaten)}

\end{large}
 % Erkl�rung (Arbeit selbstst�ndig verfasst) - auskommentiert lassen falls LMU-Arbeit
    \thispagestyle{empty}
    \cleardoublepage
    \vspace*{2cm}

\begin{center}
    \textbf{Abstract}
\end{center}

\vspace*{1cm}

\noindent Hier steht eine kurze Zusammenfassung der Arbeit. Sie darf auf gar keinen Fall
l�nger als eine Seite sein, ca. eine drittel bis eine halbe Seite ist optimal.

 % Abstract
    \thispagestyle{empty}
    \tableofcontents % Inhaltsverzeichnis

% ---------------------------------------------------------------
\mainmatter % die eigentliche Arbeit

\chapter{Introduction}
\section{Motivation}
\section{Scope}

\chapter{Background}
\section{Graph definition}
\section{Standard Template Library}
\section{Partitioned Global Address Space}
\section{DASH Library}

\chapter{Related Work}
\section{Shared Memory}
\subsection{STINGER}
\subsection{Ligra}
\section{Distributed Memory}
\subsection{Parallel Boost Graph Library}
\subsection{STAPL Parallel Graph Library}

\chapter{Container Concept}
\section{Interface semantics}
\section{Computational constraints/assumptions}
\section{Memory Space}
\section{Index Space}
\section{Iteration Space}

\chapter{Reference Implementation}

\chapter{Case studies}
\section{Static structure}
\subsection{Graph traversal}
\subsection{Shortest path evaluation}
\section{Dynamic Structure}
\subsection{Graph partitioning}
\subsection{De Bruijn Graph construction}

\chapter{Evaluation}
\section{Micro-benchmarks}

\chapter{Conclusion}
\section{Summary}
\section{Assessment}
\section{Outlook}

% ---------------------------------------------------------------
\backmatter % ab hier keine Nummerierung mehr
    \listoffigures
    \bibliographystyle{alphadin}
    \bibliography{./Bib/effe17}

\end{document}
